\section{Expected bit time}\label{sec:bit-runtime}
Here $\widetilde O$ suppresses fixed powers of $\log(n+2)$.  We count
individual unbiased random bits and bit operations, not exact arithmetic on
unbounded integers at unit cost.  It suffices to use an unweighted matching
sampler.  All stated exponents are for the order $n$, not its encoding length.

\begin{lemma}[A conservative finite-bit matching clock]\label{lem:pmclock-v2}
For a feasible zero-one board of side $s\le n$, an exact uniform
perfect matching costs $P(n)=\widetilde O(n^{21})$ expected bit operations.
\end{lemma}
\begin{proof}
We specify the bit-accounting adapter of the finite-table correction in
\cite[Sec.~6]{GLMP}.  The clock below is deliberately conservative.  Use
the lazy single-edge Jerrum--Sinclair--Vigoda chain on perfect and
near-perfect matchings of the positive completion of the board.
From perfect, propose removing a uniform matched edge; from near-perfect,
choose a uniform vertex for the corresponding plug or swap.  Use one-half
laziness and Metropolis acceptance.  Factor-two accurate hole weights
give inverse gap at most $1000s^6$ for $s\ge2$: the cross-flow bound
is $48s^4$, and the full-flow multiplier is at most $2+4(r+1/r)$ for
$r\in[s^2/2,2s^2]$, the near-to-perfect stationary mass ratio
(\href{https://www.math.cmu.edu/~af1p/Teaching/Markov_Chain_Mixing/Papers/permanent.pdf}{Jerrum--Sinclair--Vigoda,
Lemmas 4.1, 4.4}). Start each independent chain at a deterministically
found real matching $x_0$ of stationary mass at least
$[(1+2s^2)s!]^{-1}$ on good weights; side zero or one is immediate.

Let $u=C(s+2)^4(1+\lceil\log_2(s+2)\rceil)^2$ rounded up,
$\delta=2^{-u}$, with a sufficiently large fixed $C$.
At $R=O(s^2\log(s+2))$ temperatures, cool the fake edges incident to one
left vertex by the rational factor $3/4$, stopping exactly at $1/s!$.
A configuration has at most one affected edge, so adjacent ideal hole
ratios differ by at most $4/3$. At each update draw independently
$S=O(s^2(u+\log(s+2)))$ end states of
$T_0=O(s^6(s\log(s+2)+\log(s+2)))$ steps. Update each hole weight
by its old value times the ratio of the observed full to the corresponding
near-class frequency; a missing denominator or full class fails the attempt.
The initial weights are $s$.  With sufficiently large absolute coefficients,
the reversible $L^2$ mixing bound from $x_0$ gives TV error less than
$1/(1000s^2)$ per preliminary end state on good weights.  Each class
has stationary probability at least $1/[2(1+2s^2)]$ there. Independent
binomial concentration at fixed relative accuracy (say $1/50$) and a
union bound over all $R(1+s^2)$ classes give failure at most $2^{-20u}$
for the complete calibration. Indeed the ideal update is exactly
$w\pi(\mathrm{full})/\pi(\mathrm{hole})$; empirical accuracy within
a multiplicative factor $6/5$, followed by cooling, maintains the
factor-two condition.  This argument is inductive and does not multiply
per-temperature estimation errors.

At the end of this attempt draw one fresh specified chain of
$T_1=O(s^6(s\log(s+2)+u))$ steps.  On good weights the TV
error is at most $2^{-20u}$, and conditional on a real perfect endpoint
the \emph{stationary} law is exactly uniform: real matchings all have
weight one. The stationary probability of real endpoints is at least
$1/[2(1+2s^2)]$, since all fake perfect matchings together have weight
at most $s!/(s!)=1$. Repeat \emph{independent entire attempts} until
success, at most $J=O(s^2u)$ times, returning $x_0$ if exhausted.
Success has probability $\Omega(s^{-2})$, hence the expected number of
attempts is $O(s^2)$, not $J$. Bad calibration and finite mixing
together perturb the conditional success law by at most $O(s^2 2^{-20u})$.

Round token probabilities by dividing a uniform $b$-bit integer into
$4s$ almost equal bins, and round each binary acceptance probability down
to a multiple of $2^{-b}$.  Take $b=64u+O(\log(2+J(RST_0+T_1)s))$ with
sufficient padding.  Coupling all draws across the \emph{whole capped}
horizon costs at most $2^{-60u}$ in TV.  Adjusting absolute constants, the
resulting finite-output law $Q$ has TV error at most $2^{-8u}$ from the
uniform law $\pi$. Since every valid atom has $\pi(x)\ge1/s!$,
$\pi-(1-\delta)Q$ is nonnegative.  With probability $1-\delta$ use
the capped sampler; otherwise compute and sample the correction
$[\pi-(1-\delta)Q]/\delta$.

For completeness this last law is attainable at negligible \emph{expected}
cost. A single hole weight is a product of at most $R$ frequency ratios
with integers at most $S$. There are at most $(S+1)^{2Rs^2}$ weight
vectors. For each phase and vector, compute one end-state distribution
by time-layer dynamic programming on at most $(1+s^2)s!$ states.
Enumerate the histograms of $S$ independent end states on $1+s^2$
classes, at most $(S+1)^{1+s^2}$ possibilities, using the multinomial
law; this updates the vector or fails.  Compute the final end-state
distribution by the same dynamic program; apply the capped geometric sum
to the single-attempt distribution.  Enumerate legal matchings for $\pi$
only on the correction branch. Every partial law at a fixed clock time
has a common dyadic denominator of \emph{polynomial bit length}: pad
the entire attempt/retry horizon to its deterministic polynomial maximum,
with $b$ bits for a token and $b$ for acceptance at each transition.
Histogram multinomial computations likewise admit the common denominator
for $S$ padded samples.  This avoids taking coprime denominators over an
exponential number of random histories. The table work and scans, including
schoolbook arithmetic and a final exact cumulative rational selection, cost
$\exp(O(s^4\log^2(s+2)))$ bit operations in expectation. The asymptotic
coefficient in the log-width bound can be chosen independently of $C$:
$R=O(s^2\log s)$ and $\log S=O(\log s)+O(\log C)$, where the
latter additive constant contributes lower order as $s\to\infty$.
Any fixed power of the table width for straightforward scans costs the
same order in its logarithm. Take $C$ large enough that its prefactor
$2^{-u}$ pays for this whole branch.

On the fast branch, each of $O(1)$ changed edge and hole factors has
$O(R\log(S+1)+s\log(s+2))=\widetilde O(s^2)$ bits. Local probability
ratios therefore take $\widetilde O(s^4)$ schoolbook bit work.
Accessing a whole table of $s^2$ hole entries sequentially stays within
this bound. Do not construct the potentially long $b$-bit rounded
transition probability at every step. For a rational $a/d\in(0,1)$
reveal its binary digits by doubling and subtracting its remainder;
compare with fresh fair bits, stopping on the first discrepancy or rejecting
on a tie of $b$ digits. The expected number of digits used is at most two.
For $4s$ nearly uniform token bins, a dyadic prefix of length $k$ is
ambiguous only if it meets a boundary, probability $O(s2^{-k})$; the
expected prefix cost is polylogarithmic. These lazy operations have exactly
the stipulated finite dyadic law. A $2^{-u}$ correction coin likewise
short-circuits.  Thus $RST_0+T_1=\widetilde O(s^{15})$ transitions per
attempt, $O(s^2)$ expected attempts, and $\widetilde O(s^4)$ expected
bit work per transition give the asserted exponent. This is a fair-bit
adapter, not a relabeling of the arithmetic complexity in the source.
\end{proof}

\begin{proposition}[Full accounting]\label{prop:runtime}
The total expected fair-bit work, including all proposals and their
rejections and the one-time preparation, is $\widetilde O(n^{25})$.
\end{proposition}
\begin{proof}
We detail the interfaces that could otherwise conceal a further power
of $n$.  There are $O(n)$ feasible ordinary matching factors per history,
each of side at most $n$. For each designated fill of the balanced
lookahead block, use the free row and column sets of each color $k$,
with $v_k$ the size of its own pin set and $s_k=n-v_k\ge n/2$.
In the current free matching $F$ of size $s_c$, the number of assigned
edges free for the later color $k$ is $s_c-2v_k+h_k(F)$, where $h_k(F)$
counts assignments from the owner-$k$ domains to the owner-$k$ images.
The sets are disjoint and $\sum_{k>c}h_k(F)\le n$. Thus the weight is a
known rational times
\[
 \psi(F)=\prod_{k>c}(1-1/s_k)^{h_k(F)}\in[e^{-4},1]
\]
for large $n$. Rejection from an unweighted uniform matching has constant
expected cost. Its weighted normalizer is its known rational times
$\per(B)\E\psi(F)$. An independent exact $\E\psi(F)$-coin uses one
unweighted matching and one rational flip. The unbalanced fallback is
unweighted.  For each of the at most two cubic factors, write the capped
mark mass as $A\per(T)\theta_T$ with $A=n$,
$\theta_T=\E_M\min(1,2^{\#\text{cycles}(T-M)}/A)\ge1/A$.
Here too one unweighted matching gives each input coin, and a cubic
coloring draw costs at most $A$ such samples in expectation.

In Lemma~\ref{gate:binary-v2}, encode the anchor's next column by binary
splits of its current row. Conditional on all prior revelations, the
expected reciprocal of the selected branch probability is exactly the
number of positive children, at most two.  Use the deterministic upper
bound $N=O(n^2\log(n+2))$ on total sites, including the additional
scalar coins, and $m=N+2$ repetitions for the negative-binomial
estimators. Positive and reciprocal factors both use at most
$(m+1)/p$ expected probes at a site of probability $p$, also under the
specified size-biased simulation. The $O(n)$ ordinary lookahead scalar
coins have $1/p\le e^4$; the $O(1)$ cubic ones have $1/p\le n$.
Even those latter sites cost only $O(mn)$ total. The constant-mean
two-sided sigmoid race and the forward/reverse draws themselves are
therefore covered by $\widetilde O(n^4)$ unweighted matching calls per
soft gate, not $O(n^5)$ and not one weighted-graph call on a bigger board.

Here is generous accounting for everything else. In the fixed $3\times3$
cubic contingency table it suffices to index four entries and the layer
$t\le n$, since row and column sums are $t$; at most $O(n^5)$ states,
six predecessor transitions per state, and $O(n\log n)$-bit coefficients
before combining their squares and factorials. Even a full sequential
table scan per operation is below $O(n^{13})$ bit work. Each ordinary
marked-block composition table takes $O(bn^2)$ convolutions with
$O(bn\log n)$-bit entries after clearing by $n!$ at each step; the
physical block sizes have $\sum b\le n$. Fixed-length polygon and pair
cycle-index tables need only polynomially many factorials with
$O(n\log n)$ bits. Masks, feasibility via augmenting paths, row order
and shallow binary restrictions can all be scanned/constructed in
$\widetilde O(n^5)$ work outside individual probes. For rational
lookahead factors a crude bound on the aggregate known numerator and
denominator length is $O(n^4\log n)$: there are at most $n^2$ color
pairs, with exponents $O(n)$ on $O(\log n)$-bit factors. This includes
the factors on the realized matching. The logistic threshold carries
$O(n^2\log n)$ bits (and multiplication/division by the level index's
dyadic factor changes this by $O(\log n)$). Polynomially many applications
of exact rational arithmetic, with clearing of denominators as above,
fit comfortably in $O(n^{24})$ bit work. In particular, at most
$O(n^2)$ such lookahead factors can be evaluated by $O(\log n)$ repeated
squarings each. At a probe, evaluating a conditional scalar coin
directly costs at most $O(n^7)$ bit work, including a complete mask scan,
forming its exact rational, and sampling it; even this can be charged to
every probe without competing with $P(n)$.

The unbounded waiting counters are not presumed to fit in a fixed word.
Every positive binary branch has $p\ge1/n!$, and all added scalar
sites obey the bounds just given. For a negative-binomial waiting time
$V$ to $k\le m+1$ successes,
$\E[V\log(1+V)\mid p]\le(k/p)\log(1+(k+1)/p)$ by size-biased
Jensen. Uniformly its squared log length has expectation
$O((1+n\log n+\log m)^2)$ by Markov's inequality for $V$ at
exponential multiples of its mean. Consequently the $N$ unrestricted
counter factors' total bit length has second moment
$O(N^2(1+n\log n+\log m)^2)$. Multiply them without reducing and
combine with the deterministic prefactor at the end. These counter
updates, products and exact rational comparison are within the stated
$O(n^{24})$ allowance. Thus one forward or reverse soft unit has
expected cost $W(n)=\widetilde O(n^4P(n)+n^{24})$ uniformly over
inputs and dyadic levels.

Finally, the proved reverse mean bound and the incidence count lower
give $H/|\cL_n|=O(1)$ for a sufficiently large fixed multiple of the
rounded factorial scale. The explicit polynomial log-range bound $M$
in the multiscale flattening is sufficient: it contributes only
$O(\log^2(2M))$, not $M$ or $M^2$, to the expected work.
A polylogarithmically diverging safety padding of the scale also
leaves the soft exponent unchanged. The scale itself can be computed
once by factorials and outward log/exp intervals (additive error less
than one in the log); its $O(n^2\log n)$-bit length, and the
constant-degree polygon/contingency preprocessing, are below the same
overall bound. Substituting Lemma~\ref{lem:pmclock-v2} gives
$\widetilde O(n^{4+21})$. The finite small-order part can be handled
by drawing uniformly from all independent permutations and rejecting
non-Latin arrays; it does not alter the asymptotic statement. No
literal bound $O(n^{25})$, numerical crossover or worst-case
polynomial stopping time is implied by the displayed soft bound.
\end{proof}
